% !TEX root = ../thesis.tex

微小RNA（miRNA）是关键的转录后调控因子，通过调控高达60\%的蛋白质编码基因来维持细胞稳态并影响疾病进程。尽管miRNA在生物学中占据核心地位，但当前的miRNA研究仍受限于碎片化和静态的视角：上游转录调控注释不完整（起源盲区）、网络分析过度依赖单一分子的表达量而忽视调控关系（节点中心偏差）、以及缺乏对表观转录组氧化修饰的系统性审视（动态修饰缺失）。这种局限性导致了研究结果的可重复性差、机制洞见有限，以及在癌症和心血管疾病等复杂病理中的临床转化受阻。为克服这些挑战，本论文构建了一个多维度的系统性框架，旨在将miRNA研究从碎片化、静态、以单一分子为中心的现状，转向一个整体、动态、以网络为导向的新范式，从而实现基于miRNA介导的调控网络对疾病进展的精准预测和机制建模。

本研究主要包括三个递进的层面：首先，在转录起源层面，我们解决了上游调控信息缺失的难题。开发了miRStart 2.0，一个整合超过4,500个多模态数据集的深度学习平台，用于精确识别miRNA转录起始位点。该平台在与实验验证数据的对比中显示出高达93\%的准确性，成功注释了28,828个miRNA转录起始位点，与600万个组织特异性TF-miRNA调控互作，建立了一个全面的上游转录调控图谱，为后续的网络分析奠定了坚实的数据基础。其次，在网络架构层面，我们证明了网络中节点间的调控关系同样承载着关键的生物学信息。通过两个案例研究，我们展示了基于边的分析范式与优势：在三阴性乳腺癌中，基于miRNA-靶标互作构建的预后模型显著优于仅基于表达量的方法，证明了调控关系的临床稳健性；在泛疾病分析中，通过解析枢纽miR-483的网络拓扑结构，揭示了其功能二重性源于上下文依赖的网络连接模式。最后，在动态修饰层面，我们揭示了氧化修饰驱动的网络重编程效应，突破了序列静态的传统认知。我们开发了一套可从标准miRNA-seq数据中识别氧化足迹的计算流程，并引入``重编程活性评分''来定量评估氧化修饰对网络的重塑程度。应用该流程，我们在肝癌中识别出与代谢重编程密切相关的阶段特异性生物标志物，揭示了氧化应激下网络动态重连的致病机制。

综上所述，本论文通过将转录起源、网络拓扑和表观转录组动态有机结合，建立了一个系统性的miRNA调控解析框架。该框架不仅解决了miRNA研究的孤立性问题，更提升了我们预测应激诱导下疾病代谢行为的能力，为发现新型生物标志物和开发精准干预策略提供了强有力的理论与工具支持。
